\section{Introduction}
\label{sec:intro}

Robotic manipulation depends on visual reasoning that is more specific than generic image understanding. A robot must infer which object is reachable, whether a grasp is stable, which point is closer than the other in 3D, and which direction to move the end-effector next to reach a desired pose. Recent embodied and vision-language-action systems show that pretrained VLMs can transfer semantic knowledge into robot planning and control~\cite{driess2023palme,brohan2023rt2,oneill2023openx,kim2024openvla}. At the same time, spatial reasoning and grounded affordance understanding remain known weak points for VLMs, motivating targeted benchmarks instead of only broad multimodal leaderboards~\cite{liu2022vsr,chen2024spatialvlm,cheng2024spatialrgpt}.

Robo2VLM is a recent robotics VQA benchmark generated from real robot trajectories, with questions grounded in robot state, end-effector pose, force, task phase, and scene information~\cite{chen2025robo2vlm}. It is valuable because it transforms manipulation trajectories into multiple-choice VQA items. However, a broad Robo2VLM score can mix several abilities: geometric reasoning, goal-state matching, task completion recognition, obstacle reasoning, and answer-template priors. This project therefore asks a narrower question: when Robo2VLM questions are filtered for spatial reasoning and affordance understanding, do current VLMs actually solve the embodied reasoning problem, or do their scores hide shortcut behavior?

\textbf{Contributions.}
This project contributes an evaluation layer over Robo2VLM that reveals critical flaws in aggregate benchmark interpretation and in the tested VLMs:
\begin{itemize}[leftmargin=*,noitemsep,topsep=1pt]
    \item a curated Robo2VLM test split into spatial reasoning, affordance understanding, and discarded out-of-scope questions, with deterministic subcategories for depth/distance ordering, directional orientation, cross-view correspondence, object blockage/reachability, and grasp stability;
    \item a config-driven vLLM evaluation harness with zero-shot and chain-of-thought prompts, explicit option-letter scoring, RGB and RGB-depth image modes, and image-control hooks;
    \item an empirical comparison of Qwen2.5-VL-3B, DeepSeek-VL2-tiny, and poster-reported compact baselines on spatial and affordance subsets;
    \item visual-grounding sanity probes that test whether explanations remain tied to the image under correct-answer, wrong-answer, blank-image, and shuffled-image conditions.
\end{itemize}
The main finding is that spatial scores remain close to five-way chance, affordance scores are much stronger but largely driven by object blockage/reachability templates, and chain-of-thought can actively harm performance.
